\documentclass[12pt,a4paper]{article}
\usepackage[utf8]{inputenc}
\usepackage{amsmath,amsfonts,amssymb}
\usepackage{listings}
\usepackage{xcolor}
\usepackage{geometry}
\usepackage{hyperref}
\usepackage{fancyhdr}

\geometry{margin=1in}
\pagestyle{fancy}
\fancyhf{}
\rhead{Python Dictionaries - Complete Guide}
\lfoot{Python for Data Science}
\rfoot{\thepage}

\lstset{
    language=Python,
    basicstyle=\ttfamily\small,
    keywordstyle=\color{blue}\bfseries,
    commentstyle=\color{green!60!black},
    stringstyle=\color{red},
    numbers=left,
    numberstyle=\tiny\color{gray},
    stepnumber=1,
    numbersep=8pt,
    showstringspaces=false,
    breaklines=true,
    frame=single,
    backgroundcolor=\color{gray!10}
}

\title{\textbf{Python Dictionaries - Complete Guide}}
\author{Python for Data Science - Beginner's Level}
\date{\today}

\begin{document}

\maketitle
\tableofcontents
\newpage

\section{Overview}

Dictionaries are \textbf{unordered} (Python 3.7+ maintains insertion order), \textbf{mutable}, and store data in \textbf{key-value pairs}. Keys must be \textbf{unique} and \textbf{immutable}, while values can be any data type.

\section{Dictionary Creation}

\subsection{Basic Creation}

\begin{lstlisting}
# Empty dictionary
empty_dict = {}
empty_dict = dict()

# Dictionary with initial values
student = {
    "name": "Alice",
    "age": 20,
    "grade": "A",
    "subjects": ["Math", "Physics", "Chemistry"]
}

# Mixed data types
mixed_dict = {
    1: "one",
    "two": 2,
    3.0: "three point zero",
    (4, 5): "tuple key"
}
\end{lstlisting}

\subsection{Using dict() Constructor}

\begin{lstlisting}
# From keyword arguments
person = dict(name="John", age=30, city="New York")

# From list of tuples
pairs = [("a", 1), ("b", 2), ("c", 3)]
from_pairs = dict(pairs)

# From zip
keys = ["name", "age", "city"]
values = ["Alice", 25, "Boston"]
from_zip = dict(zip(keys, values))
\end{lstlisting}

\subsection{Dictionary Comprehension}

\begin{lstlisting}
# Basic comprehension
squares = {x: x**2 for x in range(6)}

# With condition
even_squares = {x: x**2 for x in range(10) if x % 2 == 0}

# From existing dictionary
original = {"a": 1, "b": 2, "c": 3}
doubled = {k: v*2 for k, v in original.items()}
\end{lstlisting}

\section{Accessing Elements}

\subsection{Basic Access}

\begin{lstlisting}
student = {"name": "Alice", "age": 20, "grade": "A"}

# Direct access (raises KeyError if key doesn't exist)
print(student["name"])  # "Alice"
print(student["age"])   # 20

# Safe access with get()
print(student.get("name"))     # "Alice"
print(student.get("height"))   # None
print(student.get("height", "Unknown"))  # "Unknown"
\end{lstlisting}

\subsection{Advanced Access Methods}

\begin{lstlisting}
student = {"name": "Alice", "age": 20, "grade": "A"}

# Check if key exists
print("name" in student)        # True
print("height" in student)      # False

# Get all keys, values, items
print(student.keys())    # dict_keys(['name', 'age', 'grade'])
print(student.values())  # dict_values(['Alice', 20, 'A'])
print(student.items())   # dict_items([('name', 'Alice'), ...])
\end{lstlisting}

\section{Modifying Dictionaries}

\subsection{Adding and Updating Elements}

\begin{lstlisting}
student = {"name": "Alice", "age": 20}

# Add new key-value pair
student["grade"] = "A"
student["subjects"] = ["Math", "Physics"]

# Update existing value
student["age"] = 21

# Update multiple values
student.update({"grade": "A+", "city": "Boston"})

# Update from another dictionary
additional_info = {"phone": "123-456-7890", "email": "alice@email.com"}
student.update(additional_info)
\end{lstlisting}

\subsection{Removing Elements}

\begin{lstlisting}
student = {
    "name": "Alice", 
    "age": 21, 
    "grade": "A+", 
    "city": "Boston",
    "phone": "123-456-7890"
}

# pop() - Remove and return value
age = student.pop("age")

# pop() with default value
height = student.pop("height", "Not specified")

# popitem() - Remove and return last inserted item
last_item = student.popitem()

# del statement
del student["city"]

# clear() - Remove all elements
student.clear()
\end{lstlisting}

\section{Dictionary Methods}

\subsection{Comprehensive Method Overview}

\begin{lstlisting}
data = {"a": 1, "b": 2, "c": 3}

# copy() - Shallow copy
data_copy = data.copy()

# fromkeys() - Create dictionary with same value for all keys
keys = ["x", "y", "z"]
default_dict = dict.fromkeys(keys, 0)

# setdefault() - Get value or set default if key doesn't exist
data.setdefault("d", 4)  # Adds "d": 4
existing = data.setdefault("a", 100)  # Returns existing value
\end{lstlisting}

\section{Iterating Through Dictionaries}

\subsection{Basic Iteration}

\begin{lstlisting}
student = {"name": "Alice", "age": 21, "grade": "A+"}

# Iterate over keys (default behavior)
for key in student:
    print(f"{key}: {student[key]}")

# Explicitly iterate over keys
for key in student.keys():
    print(f"Key: {key}")

# Iterate over values
for value in student.values():
    print(f"Value: {value}")

# Iterate over key-value pairs
for key, value in student.items():
    print(f"{key}: {value}")
\end{lstlisting}

\section{Advanced Dictionary Operations}

\subsection{Nested Dictionaries}

\begin{lstlisting}
students = {
    "alice": {
        "age": 20,
        "grades": {"math": 95, "physics": 87, "chemistry": 92},
        "contact": {"email": "alice@email.com", "phone": "123-456-7890"}
    },
    "bob": {
        "age": 22,
        "grades": {"math": 78, "physics": 85, "chemistry": 80},
        "contact": {"email": "bob@email.com", "phone": "098-765-4321"}
    }
}

# Accessing nested values
print(students["alice"]["age"])  # 20
print(students["alice"]["grades"]["math"])  # 95
\end{lstlisting}

\subsection{Dictionary as Counter}

\begin{lstlisting}
# Manual counting
text = "hello world"
char_count = {}
for char in text:
    char_count[char] = char_count.get(char, 0) + 1

# Using collections.Counter (recommended)
from collections import Counter
char_counter = Counter(text)
\end{lstlisting}

\section{Dictionary Sorting and Filtering}

\subsection{Sorting Dictionaries}

\begin{lstlisting}
data = {"banana": 3, "apple": 5, "cherry": 1, "date": 4}

# Sort by keys
sorted_by_keys = dict(sorted(data.items()))

# Sort by values
sorted_by_values = dict(sorted(data.items(), key=lambda x: x[1]))

# Sort by values (descending)
sorted_desc = dict(sorted(data.items(), key=lambda x: x[1], reverse=True))
\end{lstlisting}

\subsection{Filtering Dictionaries}

\begin{lstlisting}
data = {"apple": 5, "banana": 3, "cherry": 8, "date": 2}

# Filter by value
high_values = {k: v for k, v in data.items() if v > 4}

# Filter by key
long_names = {k: v for k, v in data.items() if len(k) > 5}

# Filter using filter() function
even_values = dict(filter(lambda x: x[1] % 2 == 0, data.items()))
\end{lstlisting}

\section{Performance Considerations}

\begin{table}[h]
\centering
\begin{tabular}{|l|c|c|l|}
\hline
\textbf{Operation} & \textbf{Average} & \textbf{Worst} & \textbf{Notes} \\
\hline
Access/Search & O(1) & O(n) & Hash collision \\
Insert & O(1) & O(n) & Hash collision \\
Delete & O(1) & O(n) & Hash collision \\
Iteration & O(n) & O(n) & All elements \\
\hline
\end{tabular}
\caption{Dictionary Operations Time Complexity}
\end{table}

\section{Common Use Cases}

\subsection{Configuration Management}

\begin{lstlisting}
config = {
    "database": {
        "host": "localhost",
        "port": 5432,
        "name": "myapp"
    },
    "api": {
        "base_url": "https://api.example.com",
        "timeout": 30
    }
}
\end{lstlisting}

\subsection{Data Transformation}

\begin{lstlisting}
# Transform list of records to dictionary
students_list = [
    {"id": 1, "name": "Alice", "grade": "A"},
    {"id": 2, "name": "Bob", "grade": "B"}
]

# Index by ID
students_by_id = {student["id"]: student for student in students_list}
\end{lstlisting}

\section{Common Pitfalls}

\subsection{Mutable Default Arguments}

\begin{lstlisting}
# Wrong way
def add_item(item, target_dict={}):  # Dangerous!
    target_dict[len(target_dict)] = item
    return target_dict

# Right way
def add_item(item, target_dict=None):
    if target_dict is None:
        target_dict = {}
    target_dict[len(target_dict)] = item
    return target_dict
\end{lstlisting}

\subsection{Key Type Restrictions}

\begin{lstlisting}
# Valid keys (immutable types)
valid_dict = {
    "string": 1,
    42: 2,
    (1, 2): 3,
    frozenset([1, 2, 3]): 4
}

# Invalid keys (mutable types)
# invalid_dict = {
#     [1, 2, 3]: "list",      # TypeError
#     {"a": 1}: "dict",       # TypeError
# }
\end{lstlisting}

\section{Summary}

Dictionaries are powerful Python data structures that provide:
\begin{itemize}
    \item \textbf{Key-value mapping}: Efficient data organization
    \item \textbf{Fast lookup}: O(1) average case access time
    \item \textbf{Flexibility}: Support for various data types
    \item \textbf{Mutability}: Can be modified after creation
    \item \textbf{Rich methods}: Comprehensive built-in functionality
    \item \textbf{Memory efficiency}: Optimized hash table implementation
\end{itemize}

Master dictionaries to handle complex data relationships and build efficient applications!

\end{document}